%=============================================================
%  SEGMENT 5 — Back Matter: Appendices, Bibliography, Index
%  Book: وعده یا خدعه؟ — بررسی انتقادی
%  Codename: khomeini-promise-critical-study
%  Requires: Segments 0–4 compiled; XeLaTeX ×3 (for index)
%=============================================================

%%%%%%%%%%%%%%%%%%%%%%%%%%%%%%%%%%%%%%%%%%%%%%%%%%%%%%%%%%%%
%%  APPENDICES
%%%%%%%%%%%%%%%%%%%%%%%%%%%%%%%%%%%%%%%%%%%%%%%%%%%%%%%%%%%%
\appendix
\renewcommand{\chaptername}{پیوست}

%%%%%%%%%%%%%%%%%%%%%%%%%%%%%%%%%%%%%%%%%%%%%%%%%%%%%%%%%%%%
%%  APPENDIX A — گاه‌شمار کامل
%%%%%%%%%%%%%%%%%%%%%%%%%%%%%%%%%%%%%%%%%%%%%%%%%%%%%%%%%%%%
\chapter{گاه‌شمار تفصیلی: از نجف تا درگذشت خمینی}
\label{app:timeline}

\begin{infoBox}[دربارهٔ این پیوست]
این پیوست گاه‌شمار تفصیلیِ رویدادهای کلیدی
از تبعید خمینی به نجف (۱۳۴۳) تا درگذشت او (۱۳۶۸)
را ارائه می‌دهد.
هر رویداد با تز یا تزهای مرتبط رنگ‌کدگذاری شده است.
\end{infoBox}

{%
\small
\begin{longtable}{%
  >{\centering\arraybackslash}m{2cm}
  >{\centering\arraybackslash}m{2.5cm}
  >{\arraybackslash}m{6cm}
  >{\centering\arraybackslash}m{2.2cm}}
\caption{گاه‌شمار تفصیلی}
\label{tab:full-timeline} \\
\toprule
\textbf{سال شمسی}
  & \textbf{سال میلادی}
  & \textbf{رویداد}
  & \textbf{تز مرتبط} \\
\midrule
\endfirsthead
\caption[]{(ادامه)} \\
\toprule
\textbf{سال شمسی}
  & \textbf{سال میلادی}
  & \textbf{رویداد}
  & \textbf{تز مرتبط} \\
\midrule
\endhead
\midrule
\multicolumn{4}{r}{\footnotesize \textit{ادامه در صفحهٔ بعد}\,\dots} \\
\endfoot
\bottomrule
\endlastfoot

%--- pre-revolution ---
۱۳۴۳
  & ۱۹۶۴
  & تبعید خمینی به ترکیه و سپس نجف%
    \footnote{\en{Moin},
    \textit{\en{Khomeini: Life of the Ayatollah}}, 1999, pp.\,104–112.}
  & — \\
\addlinespace

۱۳۴۸
  & ۱۹۷۰
  & انتشار درس‌گفتارهای \textit{ولایت فقیه}
    (\en{Islamic Government}) در نجف%
    \footnote{\en{Martin},
    \textit{\en{Creating an Islamic State}}, 2000, pp.\,78–95.}
  & \cellcolor{cGold!15}ت‌۳ / \cellcolor{cAccent!15}ت‌۲ \\
\addlinespace

۱۳۵۶
  & ۱۹۷۸
  & آغاز اعتراضات گسترده در ایران؛
    خمینی از نجف به پاریس می‌رود (مهر ۱۳۵۷)%
    \footnote{\en{Keddie},
    \textit{\en{Modern Iran}}, 2006, pp.\,216–224.}
  & — \\
\addlinespace

۱۳۵۷ (دی–بهمن)
  & ۱۹۷۹ (ژانویه)
  & مصاحبه‌های متعدد پاریس:
    وعدهٔ عدم تصدّی مقام سیاسی%
    \footnote{صحیفهٔ امام، ج\,۵، صص\,۳۸۸–۳۹۰.}
  & \cellcolor{cGreen!15}ت‌۱ / \cellcolor{cAccent!15}ت‌۲ \\
\addlinespace

%--- revolution ---
۱۳۵۷ (۲۲ بهمن)
  & ۱۹۷۹ (۱۱ فوریه)
  & پیروزی انقلاب؛ سقوط حکومت بختیار%
    \footnote{بازرگان، \textit{انقلاب ایران در دو حرکت}، ص\,۷۲.}
  & \cellcolor{cAccent!15}ت‌۲ / \cellcolor{cConsolid!15}ت‌۴ \\
\addlinespace

۱۳۵۷ (اسفند)
  & ۱۹۷۹ (مارس)
  & رفراندوم جمهوری اسلامی (۱۲ فروردین ۱۳۵۸)%
    \footnote{\en{Milani},
    \textit{\en{The Making of Iran's Islamic Revolution}},
    1994, p.\,156.}
  & \cellcolor{cAccent!15}ت‌۲ \\
\addlinespace

۱۳۵۸ (تابستان)
  & ۱۹۷۹
  & مجلس خبرگان قانون اساسی:
    طرح ولایت فقیه مطلقه%
    \footnote{صورت مشروح مذاکرات مجلس خبرگان قانون اساسی، ۱۳۵۸.}
  & \cellcolor{cAccent!15}ت‌۲ / \cellcolor{cGold!15}ت‌۳ \\
\addlinespace

۱۳۵۸ (آبان)
  & ۱۹۷۹ (نوامبر)
  & اشغال سفارت آمریکا — آغاز بحران گروگان‌گیری%
    \footnote{\en{Bowden},
    \textit{\en{Guests of the Ayatollah}}, 2006, pp.\,1–25.}
  & \cellcolor{cAccent!15}ت‌۲ / \cellcolor{cConsolid!15}ت‌۴ \\
\addlinespace

۱۳۵۸ (آذر)
  & ۱۹۷۹ (دسامبر)
  & رفراندوم قانون اساسی — تصویب ولایت فقیه%
    \footnote{\en{Schirazi},
    \textit{\en{The Constitution of Iran}}, 1997, pp.\,22–48.}
  & \cellcolor{cGold!15}ت‌۳ \\
\addlinespace

۱۳۵۸–۱۳۵۹
  & ۱۹۷۹–۱۹۸۰
  & ترورهای گروه فرقان؛
    ترور مطهری، مفتح و دیگران%
    \footnote{\en{Abrahamian},
    \textit{\en{Radical Islam}}, 1989, pp.\,186–193.}
  & \cellcolor{cConsolid!15}ت‌۴ \\
\addlinespace

۱۳۵۹ (خرداد)
  & ۱۹۸۰ (ژوئن)
  & عزل بنی‌صدر از فرماندهی کل قوا%
    \footnote{بنی‌صدر، \textit{خیانت به امید}، صص\,۳۰۲–۳۱۵.}
  & \cellcolor{cAccent!15}ت‌۲ / \cellcolor{cConsolid!15}ت‌۴ \\
\addlinespace

۱۳۵۹ (شهریور)
  & ۱۹۸۰ (سپتامبر)
  & آغاز جنگ ایران و عراق%
    \footnote{\en{Crist},
    \textit{\en{The Twilight War}}, 2012, pp.\,42–56.}
  & \cellcolor{cGold!15}ت‌۳ / \cellcolor{cConsolid!15}ت‌۴ \\
\addlinespace

۱۳۶۰ (خرداد–تیر)
  & ۱۹۸۱
  & سرکوب مجاهدین خلق و تصفیهٔ سیاسی گسترده%
    \footnote{\en{Abrahamian},
    \textit{\en{Tortured Confessions}}, 1999, pp.\,124–158.}
  & \cellcolor{cAccent!15}ت‌۲ / \cellcolor{cConsolid!15}ت‌۴ \\
\addlinespace

۱۳۶۱–۱۳۶۵
  & ۱۹۸۲–۱۹۸۶
  & ادامهٔ جنگ پس از آزادسازی خرمشهر؛
    رد پیشنهادهای صلح%
    \footnote{خاطرات رفسنجانی، \textit{دوران دفاع مقدس}،
    ج\,۱، صص\,۱۱۲–۱۲۸.}
  & \cellcolor{cGold!15}ت‌۳ / \cellcolor{cConsolid!15}ت‌۴ \\
\addlinespace

۱۳۶۶
  & ۱۹۸۸
  & نامهٔ معروف خمینی دربارهٔ تقدم مصلحت نظام%
    \footnote{صحیفهٔ امام، ج\,۲۱.}
  & \cellcolor{cGold!15}ت‌۳ \\
\addlinespace

۱۳۶۷ (تیر)
  & ۱۹۸۸ (ژوئیه)
  & پذیرش قطعنامهٔ ۵۹۸ — «نوشیدن جام زهر»%
    \footnote{\en{Crist}, 2012, pp.\,350–362.}
  & \cellcolor{cGold!15}ت‌۳ \\
\addlinespace

۱۳۶۷ (مرداد–شهریور)
  & ۱۹۸۸ (اوت–سپتامبر)
  & کشتار زندانیان سیاسی ۱۳۶۷%
    \footnote{\en{Robertson},
    \textit{\en{Massacre of Political Prisoners in Iran, 1988}},
    2010, pp.\,54–87.}
  & \cellcolor{cAccent!15}ت‌۲ / \cellcolor{cConsolid!15}ت‌۴ \\
\addlinespace

۱۳۶۷ (فروردین ۱۳۶۸)
  & ۱۹۸۹ (مارس)
  & عزل آیت‌الله منتظری از قائم‌مقامی رهبری%
    \footnote{خاطرات منتظری (نسخهٔ اینترنتی).}
  & \cellcolor{cAccent!15}ت‌۲ \\
\addlinespace

۱۳۶۸ (۱۴ خرداد)
  & ۱۹۸۹ (۳ ژوئن)
  & درگذشت خمینی%
    \footnote{\en{Moin}, 1999, pp.\,305–312.}
  & — \\

\end{longtable}
}

%============================================================
\section{نمودار خطّ زمانیِ کامل}
\label{sec:app-a-tikz}

\begin{figure}[ht]
\centering
\begin{tikzpicture}[
    yscale=0.65,
    evt/.style={rectangle, rounded corners=3pt, draw=#1, thick,
                fill=#1!8, font=\tiny, align=right,
                text width=4.5cm, inner sep=3pt},
    yearmark/.style={circle, fill=cPrimary, text=white,
                     font=\tiny\bfseries, inner sep=2pt,
                     minimum size=0.7cm},
  ]
  % vertical timeline
  \draw[very thick, cPrimary!40] (0,0) -- (0,-20);

  % helper macro
  \newcommand{\timeevt}[4]{%
    % #1=y, #2=color, #3=year, #4=text
    \node[yearmark] at (0,#1) {#3};
    \node[evt=#2, anchor=west] at (1,#1) {#4};
  }

  \timeevt{0}{cGold}{۴۸}{انتشار \textit{ولایت فقیه} — نجف}
  \timeevt{-2}{cAccent}{۵۷}{وعده‌های پاریس — عدم حاکمیت}
  \timeevt{-4}{cAccent}{۵۷}{پیروزی انقلاب — نصب بازرگان}
  \timeevt{-6}{cAccent}{۵۸}{مجلس خبرگان — ولایت فقیه در قانون اساسی}
  \timeevt{-8}{cConsolid}{۵۸}{بحران گروگان‌گیری — حذف رقبا}
  \timeevt{-10}{cConsolid}{۵۹}{ترورهای فرقان — تحکیم نهادها}
  \timeevt{-12}{cAccent}{۶۰}{عزل بنی‌صدر — سرکوب مجاهدین}
  \timeevt{-14}{cGold}{۶۰–۶۶}{جنگ هشت‌ساله — بسیج و مصلحت نظام}
  \timeevt{-16}{cGold}{۶۶}{نامهٔ مصلحت نظام — ولایت مطلقه}
  \timeevt{-18}{cAccent}{۶۷}{کشتار ۱۳۶۷ — فتوای اعدام جمعی}
  \timeevt{-20}{cPrimary}{۶۸}{درگذشت خمینی}

  % power gradient bar
  \shade[top color=cGreen!30, bottom color=cAccent!50]
    (-1.8,0.5) rectangle (-1.2,-20.5);
  \node[font=\tiny, rotate=90, text=cPrimary] at (-2.2,-10)
    {وعده \quad$\longrightarrow$\quad حاکمیت مطلقه};
\end{tikzpicture}
\caption{خط زمانی عمودی: مسیر انباشت قدرت}
\label{fig:full-vertical-timeline}
\end{figure}

\begin{principleBox}[خلاصهٔ پیوست الف]
\begin{itemize}
\item گاه‌شمار تفصیلی ۱۸ رویداد کلیدی را از ۱۳۴۳ تا ۱۳۶۸
      فهرست کرد.
\item هر رویداد با تز یا تزهای مرتبط رنگ‌کدگذاری شد.
\item خط زمانی عمودی روند انباشت قدرت را
      به‌صورت بصری تأیید می‌کند.
\end{itemize}
\end{principleBox}


%%%%%%%%%%%%%%%%%%%%%%%%%%%%%%%%%%%%%%%%%%%%%%%%%%%%%%%%%%%%
%%  APPENDIX B — متون کلیدی
%%%%%%%%%%%%%%%%%%%%%%%%%%%%%%%%%%%%%%%%%%%%%%%%%%%%%%%%%%%%
\chapter{گزیدهٔ متون کلیدی}
\label{app:texts}

\begin{infoBox}[دربارهٔ این پیوست]
این پیوست گزیده‌ای از مهم‌ترین متون اولیه
— مصاحبه‌ها، سخنرانی‌ها، نامه‌ها و فتاوا —
را که در تحلیل‌های فصل‌های ۳–۱۶ مورد استناد قرار گرفتند،
به‌صورت اصلی (یا ترجمهٔ فارسیِ معتبر) ارائه می‌دهد.
هدف، فراهم‌سازیِ امکان ارجاع مستقیم
برای خواننده است.%
\footnote{تمامی متون از نسخه‌های رسمیِ منتشرشده نقل شده‌اند.
در مواردی که نسخهٔ اینترنتی مورد استفاده بوده، نشانیِ وب
در پانوشت ذکر شده است.}
\end{infoBox}

%============================================================
\section{مصاحبهٔ خمینی با \en{Le Monde} (دی ۱۳۵۷)}
\label{sec:app-b-lemonde}

\begin{warningBox}[متن اصلی — ترجمهٔ فارسی]
«\,من در آیندهٔ ایران هیچ نقش حکومتی نخواهم داشت.
نقش من هدایت و ارشاد است.
من یک آخوند هستم و جایگاه من حوزهٔ علمیه است.
حکومت به دست مردم و نمایندگان منتخب آن‌ها اداره خواهد شد.\,»%
\footnote{صحیفهٔ امام، ج\,۵، صص\,۳۸۸–۳۸۹.
متن فرانسوی اصلی:
\en{``Je n'exercerai aucune fonction de gouvernement.
Mon rôle est celui d'un guide\dots''},
\textit{\en{Le Monde}}, 9 janvier 1979.}
\end{warningBox}

\begin{noteBox}[تحلیل مختصر]
این متن مهم‌ترین سند وعدهٔ عدم حاکمیت است.
عبارت «هیچ نقش حکومتی» صریح و بدون ابهام است.
تز ۲ (فریب) این متن را تاکتیکی می‌خواند؛
تز ۱ (صداقت) آن را بیان واقعی نیت می‌داند؛
تز ۳ (تطور) معتقد است نیت بعداً تغییر کرد.
ر.ک.\ فصل~۶ برای تحلیل تفصیلی.
\end{noteBox}

%============================================================
\section{مصاحبهٔ خمینی با \en{The Guardian} (بهمن ۱۳۵۷)}
\label{sec:app-b-guardian}

\begin{warningBox}[متن اصلی — ترجمهٔ فارسی]
«\,من قصد ندارم رئیس‌جمهور شوم یا نخست‌وزیر.
من در قم خواهم ماند و از آنجا نظارت خواهم کرد.\,»%
\footnote{صحیفهٔ امام، ج\,۵، ص\,۴۲۳.
\en{The Guardian}, 6 November 1978 (interview conducted
in Neauphle-le-Château).}
\end{warningBox}

%============================================================
\section{درس‌گفتارهای \textit{ولایت فقیه} — نجف (۱۳۴۸)}
\label{sec:app-b-velayat}

\begin{warningBox}[گزیدهٔ متن]
«\,همان‌طور که پیامبر (ص) ولایت بر امت داشت
و حکومت تشکیل داد\dots\
فقیه عادل نیز باید حکومت تشکیل دهد\dots\
ولایت فقیه از جنس ولایت پیامبر و ائمه است.\,»%
\footnote{خمینی، \textit{ولایت فقیه (حکومت اسلامی)}،
نشر مؤسسهٔ تنظیم و نشر آثار امام خمینی، صص\,۳۵–۴۰.}
\end{warningBox}

\begin{noteBox}[تحلیل مختصر]
این متن — که نُه سال پیش از وعده‌های پاریس نوشته شده —
مستقیماً ایدهٔ حکومت فقیه را مطرح می‌کند.
تز ۲ آن را دلیل قاطع بر فریبکارانه بودن وعده‌های پاریس می‌داند.
تز ۳ استدلال می‌کند که این متن فقط مرحله‌ای از تطور بود
و خمینی ممکن است بعداً نظرش را تعدیل کرده باشد.
ر.ک.\ فصل‌های ۴ و ۸.%
\footnote{\en{Martin}, 2000, pp.\,78–95;
کدیور، \textit{نظریه‌های دولت در فقه شیعه}، صص\,۱۵۲–۱۸۰.}
\end{noteBox}

%============================================================
\section{نامهٔ خمینی دربارهٔ مصلحت نظام (دی ۱۳۶۶)}
\label{sec:app-b-maslahat}

\begin{warningBox}[گزیدهٔ متن]
«\,حکومت\dots\ از اهمّ احکام الهی است
و بر جمیع احکام فرعیهٔ الهیه تقدّم دارد\dots\
حکومت می‌تواند هر امری را — چه عبادی و چه غیرعبادی —
که جریان آن مخالف مصالح اسلام است
مادامی که چنین است جلوگیری کند.\,»%
\footnote{صحیفهٔ امام، ج\,۲۱، نامهٔ ۱۶ دی ۱۳۶۶
خطاب به آیت‌الله خامنه‌ای (رئیس‌جمهور وقت).}
\end{warningBox}

\begin{noteBox}[تحلیل مختصر]
این نامه نقطهٔ اوجِ تطور ایدئولوژیک (ت‌۳) است:
حکومت فقیه نه‌تنها فرعی نیست
بلکه بر تمام احکام دیگر اولویت دارد.
فاصلهٔ این موضع با وعدهٔ «هیچ نقش حکومتی»
بیش از آنکه صرفاً با «اضطرار» (ت‌۱) توضیح‌پذیر باشد،
نشان‌دهندهٔ تحولی بنیادین در اندیشه یا
آشکارسازیِ اندیشه‌ای پنهان (ت‌۲) است.
ر.ک.\ فصل‌های ۱۱ و ۱۵.%
\footnote{کدیور، \textit{حکومت ولایی}، صص\,۲۵۵–۲۷۰.}
\end{noteBox}

%============================================================
\section{نامهٔ منتظری به خمینی (مرداد ۱۳۶۷)}
\label{sec:app-b-montazeri}

\begin{warningBox}[گزیدهٔ متن]
«\,آیا می‌دانید در زندان‌های جمهوری اسلامی
زنان باردار اعدام می‌شوند؟\dots\
جهانیان ما را به کشتار متهم خواهند کرد\dots\
تاریخ ما را جنایتکار ثبت خواهد کرد.\,»%
\footnote{خاطرات آیت‌الله منتظری (نسخهٔ اینترنتی)،
نامهٔ ۹ مرداد ۱۳۶۷.
نیز ر.ک.\ \en{Robertson}, 2010, pp.\,75–82.}
\end{warningBox}

\begin{noteBox}[تحلیل مختصر]
این نامه مهم‌ترین سند اعتراض درون‌حکومتی
به کشتار ۱۳۶۷ است.
پاسخ خمینی — عزل منتظری از قائم‌مقامی رهبری —
نشان‌دهندهٔ اوج تمرکز قدرت فردی است (ت‌۲ + ت‌۴).
حتی نزدیک‌ترین متحد ایدئولوژیک نیز
قادر به تغییر مسیر نبود.
ر.ک.\ فصل~۱۲.
\end{noteBox}

\begin{principleBox}[خلاصهٔ پیوست ب]
\begin{itemize}
\item پنج متن کلیدی ارائه شد:
      مصاحبه‌های پاریس (\en{Le Monde} و \en{The Guardian})،
      درس‌گفتارهای ولایت فقیه (۱۳۴۸)،
      نامهٔ مصلحت نظام (۱۳۶۶)،
      و نامهٔ منتظری (۱۳۶۷).
\item هر متن با تحلیل مختصر و ارجاع به فصل مربوطه همراه شد.
\item هدف: فراهم‌سازیِ دسترسی مستقیم خوانندگان
      به شواهد اولیه.
\end{itemize}
\end{principleBox}


%%%%%%%%%%%%%%%%%%%%%%%%%%%%%%%%%%%%%%%%%%%%%%%%%%%%%%%%%%%%
%%  APPENDIX C — واژه‌نامهٔ تخصصی
%%%%%%%%%%%%%%%%%%%%%%%%%%%%%%%%%%%%%%%%%%%%%%%%%%%%%%%%%%%%
\chapter{واژه‌نامهٔ تخصصی (فارسی–انگلیسی)}
\label{app:glossary}

\begin{infoBox}[دربارهٔ این پیوست]
واژه‌نامهٔ زیر مهم‌ترین اصطلاحات تخصصی به‌کاررفته
در این کتاب را همراه با معادل انگلیسی
و ارجاع به فصل مربوطه فهرست می‌کند.
\end{infoBox}

{%
\small
\begin{longtable}{%
  >{\arraybackslash}m{3.5cm}
  >{\arraybackslash}m{4cm}
  >{\arraybackslash}m{5.5cm}}
\caption{واژه‌نامهٔ تخصصی}
\label{tab:glossary} \\
\toprule
\textbf{اصطلاح فارسی}
  & \textbf{معادل انگلیسی}
  & \textbf{توضیح / ارجاع} \\
\midrule
\endfirsthead
\caption[]{(ادامه)} \\
\toprule
\textbf{اصطلاح فارسی}
  & \textbf{معادل انگلیسی}
  & \textbf{توضیح / ارجاع} \\
\midrule
\endhead
\bottomrule
\endlastfoot

ولایت فقیه
  & \en{Guardianship of the Jurist}
  & نظریهٔ حکومت مستقیم فقیه؛ فصل‌های ۴، ۸ \\
\addlinespace

ولایت مطلقهٔ فقیه
  & \en{Absolute Guardianship of the Jurist}
  & نسخهٔ گسترش‌یافته: ولایت بر تمام احکام؛ فصل‌های ۱۱، ۱۵ \\
\addlinespace

مصلحت نظام
  & \en{Expediency of the State / Maṣlaḥa}
  & تقدم مصلحت حکومت بر احکام فقهی؛ فصل‌های ۱۱، ۱۶ \\
\addlinespace

تقیه
  & \en{Taqiyya (precautionary dissimulation)}
  & پنهان‌سازی باور واقعی در شرایط خطر؛ فصل~۴ \\
\addlinespace

ناهماهنگی شناختی
  & \en{Cognitive Dissonance}
  & تنش میان باور و رفتار؛ فصل~۱۵ \\
\addlinespace

خودفریبی
  & \en{Self-Deception}
  & باور و انکار هم‌زمان؛ فصل‌های ۱۴–۱۶ \\
\addlinespace

استدلال انگیزه‌محور
  & \en{Motivated Reasoning}
  & پردازش گزینشی شواهد؛ فصل~۱۵ \\
\addlinespace

سوگیری تأیید
  & \en{Confirmation Bias}
  & توجه بیشتر به شواهد هم‌سو با باور موجود؛ فصل~۱۵ \\
\addlinespace

تحلیل حساسیت
  & \en{Sensitivity Analysis}
  & بررسی تأثیر تغییر فروض بر نتیجه؛ فصل~۱۵ \\
\addlinespace

ارزیابی بیزی کیفی
  & \en{Qualitative Bayesian Assessment}
  & روش وزن‌دهی با پیشین و درست‌نمایی؛ فصل~۱۵ \\
\addlinespace

تز ترکیبی مرحله‌ای
  & \en{Phased Composite Thesis}
  & فرضیهٔ نهایی کتاب: ترکیب مرحله‌ای تزها؛ فصل‌های ۱۴–۱۶ \\
\addlinespace

پارسیمونی
  & \en{Parsimony (Occam's Razor)}
  & ترجیح تبیین ساده‌تر؛ فصل~۱۴ \\
\addlinespace

ماتریس سازگاری
  & \en{Compatibility Matrix}
  & جدول سازگاری منطقی تزها؛ فصل~۱۴ \\
\addlinespace

سوگیری بقا
  & \en{Survivorship Bias}
  & حفظ نامتوازن شواهد؛ فصل~۱۵ \\
\addlinespace

نظریهٔ ذهن
  & \en{Theory of Mind}
  & توانایی درک حالات ذهنی دیگران؛ فصل~۱۴ \\

\end{longtable}
}

\begin{principleBox}[خلاصهٔ پیوست ج]
\begin{itemize}
\item ۱۵ اصطلاح تخصصی فارسی–انگلیسی فهرست شد.
\item هر اصطلاح با ارجاع به فصل مربوطه همراه است.
\item این واژه‌نامه ابزار مرجع سریع برای خوانندگان است.
\end{itemize}
\end{principleBox}


%%%%%%%%%%%%%%%%%%%%%%%%%%%%%%%%%%%%%%%%%%%%%%%%%%%%%%%%%%%%
%%  APPENDIX D — روش‌شناسی تفصیلی
%%%%%%%%%%%%%%%%%%%%%%%%%%%%%%%%%%%%%%%%%%%%%%%%%%%%%%%%%%%%
\chapter{تشریح روش‌شناسی پژوهش}
\label{app:method}

\begin{infoBox}[دربارهٔ این پیوست]
این پیوست روش‌شناسیِ به‌کاررفته در کتاب را
به‌صورت تفصیلی شرح می‌دهد:
معیارهای انتخاب منابع، مقیاس ارزیابی،
محدودیت‌ها و مفروضات.
\end{infoBox}

%============================================================
\section{معیارهای انتخاب منابع}
\label{sec:app-d-sources}

\begin{table}[ht]
\centering
\caption{سلسله‌مراتب منابع}
\label{tab:source-hierarchy}
\begin{tabular}{>{\centering\arraybackslash}m{1.5cm}
                >{\arraybackslash}m{4cm}
                >{\arraybackslash}m{6.5cm}}
\toprule
\textbf{رتبه} & \textbf{نوع منبع} & \textbf{مثال} \\
\midrule
۱ & اسناد دست‌اول رسمی
  & صحیفهٔ امام، مصوبات مجلس خبرگان، اسناد لانهٔ جاسوسی \\
\addlinespace
۲ & خاطرات دست‌اول
  & رفسنجانی، منتظری، یزدی، بازرگان، بنی‌صدر \\
\addlinespace
۳ & پژوهش‌های آکادمیک (فارسی و انگلیسی)
  & \en{Abrahamian, Bakhash, Milani, Keddie};
    کدیور، بشیریه، سروش \\
\addlinespace
۴ & تحلیل‌های روزنامه‌نگاری و مستند
  & \en{Bowden, Crist}; مقالات تحلیلی \\
\bottomrule
\end{tabular}
\end{table}

\begin{noteBox}[اصل منبع‌شناختی]
در تمام فصل‌ها، تلاش شده منابع ردهٔ ۱ و ۲
بر ردهٔ ۳ و ۴ اولویت داشته باشند.
هرجا منبع دست‌اول در دسترس نبود،
حداقل دو منبع ثانویهٔ مستقل مورد استناد قرار گرفته است.%
\footnote{%
\en{George \& Bennett},
\textit{\en{Case Studies and Theory Development}}, 2005, pp.\,89–95.}
\end{noteBox}

%============================================================
\section{مقیاس ارزیابی سازگاری}
\label{sec:app-d-scale}

مقیاس سه‌درجه‌ای (فصل~۱۳) بر پایهٔ سه معیار عمل می‌کند:

\begin{enumerate}
\item \textbf{سازگاری مستقیم:}
  آیا شواهد مستقیماً پیش‌بینیِ تز را تأیید می‌کنند؟
\item \textbf{نیاز به فروض کمکی:}
  آیا تز برای توضیح رویداد نیاز به فروض اضافی دارد؟
\item \textbf{وجود شواهد مخالف:}
  آیا شواهدی وجود دارد که مستقیماً
  پیش‌بینی تز را نقض کنند؟
\end{enumerate}

\begin{figure}[ht]
\centering
\begin{tikzpicture}[
    node distance=0.8cm and 2cm,
    decision/.style={diamond, draw=cPrimary, thick, fill=cLight,
                     aspect=2.5, font=\scriptsize, align=center,
                     inner sep=2pt},
    result/.style={rectangle, rounded corners=4pt, draw=#1, thick,
                   fill=#1!15, font=\scriptsize\bfseries,
                   minimum width=2.2cm, minimum height=0.7cm},
    arr/.style={-{Stealth[length=5pt]}, thick},
  ]
  \node[decision] (Q1) {شواهد\\مستقیم؟};
  \node[decision, below left=1.5cm and 1.5cm of Q1] (Q2) {فروض\\کمکی؟};
  \node[result=cGreen, right=2.5cm of Q1] (R3) {قوی (۳)};
  \node[result=cGold, right=2.5cm of Q2] (R2) {متوسط (۲)};
  \node[result=cAccent, below=1.5cm of Q2] (R1) {ضعیف (۱)};

  \draw[arr] (Q1) -- node[above, font=\tiny]{بله} (R3);
  \draw[arr] (Q1) -- node[left, font=\tiny]{خیر} (Q2);
  \draw[arr] (Q2) -- node[above, font=\tiny]{بله و کافی} (R2);
  \draw[arr] (Q2) -- node[left, font=\tiny]{خیر یا ناکافی} (R1);
\end{tikzpicture}
\caption{فلوچارت تصمیم‌گیری برای نمره‌دهی سازگاری}
\label{fig:scoring-flowchart}
\end{figure}

%============================================================
\section{مفروضات و محدودیت‌ها}
\label{sec:app-d-limits}

\begin{enumerate}
\item \textbf{فرض عقلانیت حداقلی:}
  خمینی به‌عنوان فاعل عقلانی (نه لزوماً کاملاً عقلانی)
  فرض شده است — یعنی کنش‌هایش حداقل
  تا حدّی هدفمند و محاسبه‌شده بوده‌اند.%
  \footnote{\en{Jon Elster},
  \textit{\en{Ulysses and the Sirens}}, rev.\ ed., 1984, ch.\,1.}

\item \textbf{عدم دسترسی به آرشیو محرمانه:}
  بسیاری از اسناد داخلی حکومت ایران
  طبقه‌بندی‌شده و غیرقابل دسترسی‌اند.
  این محدودیت مهم‌ترین عامل عدم قطعیت نتایج است.

\item \textbf{سوگیری بقا:}
  شواهدی که به نفع حاکمیت بودند
  بیشتر حفظ و منتشر شده‌اند.
  شواهد تردید یا اعتراض ممکن است
  به‌طور سیستماتیک حذف شده باشند.

\item \textbf{زبان و ترجمه:}
  بخشی از مصاحبه‌های پاریس به فرانسه یا انگلیسی بود
  و ترجمهٔ فارسی ممکن است ظرایفی را از دست داده باشد.%
  \footnote{ر.ک.\ فصل~۶ برای بحث دربارهٔ نسخه‌های مختلف
  مصاحبهٔ خمینی با \en{Le Monde}.}
\end{enumerate}

\begin{principleBox}[خلاصهٔ پیوست د]
\begin{itemize}
\item سلسله‌مراتب چهارسطحی منابع تعریف شد.
\item مقیاس سه‌درجه‌ای سازگاری
      با فلوچارت تصمیم‌گیری تشریح شد.
\item چهار مفروض و محدودیت اصلی تصریح شدند.
\end{itemize}
\end{principleBox}


%%%%%%%%%%%%%%%%%%%%%%%%%%%%%%%%%%%%%%%%%%%%%%%%%%%%%%%%%%%%
%%  BIBLIOGRAPHY — کتاب‌نامه
%%%%%%%%%%%%%%%%%%%%%%%%%%%%%%%%%%%%%%%%%%%%%%%%%%%%%%%%%%%%

% If using BibLaTeX (recommended for this project):
% \printbibliography[heading=bibintoc, title={کتاب‌نامه}]
%
% Below: manual bibliography as fallback
% (replace with .bib file for production)

\cleardoublepage
\phantomsection
\addcontentsline{toc}{chapter}{کتاب‌نامه}
\chapter*{کتاب‌نامه}
\label{ch:bibliography}

\begin{infoBox}[دربارهٔ کتاب‌نامه]
کتاب‌نامه به دو بخش تقسیم شده است:
(الف) منابع فارسی و (ب) منابع انگلیسی.
ترتیب هر بخش الفبایی بر اساس نام خانوادگی مؤلف است.
\end{infoBox}

%============================================================
\section*{الف) منابع فارسی}

\begin{enumerate}[label={[\arabic*]}, leftmargin=2cm]

\item
بازرگان، مهدی.
\textit{انقلاب ایران در دو حرکت}.
تهران، ۱۳۶۳.

\item
بشیریه، حسین.
\textit{دیباچه‌ای بر جامعه‌شناسی سیاسی ایران (دورهٔ جمهوری اسلامی)}.
تهران: نگاه معاصر، ۱۳۸۱.

\item
بنی‌صدر، ابوالحسن.
\textit{خیانت به امید}.
پاریس، ۱۳۶۱.

\item
حجاریان، سعید.
\textit{از شاهد قدسی تا شاهد بازاری}.
تهران: طرح نو، ۱۳۸۰.

\item
خمینی، روح‌الله.
\textit{صحیفهٔ امام}، ۲۲ جلد.
تهران: مؤسسهٔ تنظیم و نشر آثار امام خمینی، ۱۳۷۸.

\item
خمینی، روح‌الله.
\textit{ولایت فقیه (حکومت اسلامی)}.
نجف، ۱۳۴۸ ش. [تقریرات درس].

\item
رفسنجانی، اکبر هاشمی.
\textit{خاطرات} (چندین جلد).
تهران: دفتر نشر معارف انقلاب، سال‌های مختلف.

\item
سروش، عبدالکریم.
\textit{فربه‌تر از ایدئولوژی}.
تهران: صراط، ۱۳۷۲.

\item
\textit{صورت مشروح مذاکرات مجلس خبرگان قانون اساسی}.
تهران: ادارهٔ کل امور فرهنگی و روابط عمومی مجلس شورای اسلامی، ۱۳۶۴.

\item
کدیور، محسن.
\textit{حکومت ولایی}.
تهران: نشر نی، ۱۳۷۷.

\item
کدیور، محسن.
\textit{نظریه‌های دولت در فقه شیعه}.
تهران: نشر نی، ۱۳۷۶.

\item
گنجی، اکبر.
\textit{تلقی فاشیستی از دین و حکومت}.
تهران: طرح نو، ۱۳۷۹.

\item
منتظری، حسینعلی.
\textit{خاطرات} (نسخهٔ اینترنتی).
\url{https://amontazeri.com/farsi/memoir}.

\item
میلانی، عباس.
\textit{نگاهی به شاه}.
تهران: نشر اختران، ۱۳۹۲.

\item
یزدی، ابراهیم.
\textit{آخرین تلاش‌ها در آخرین روزها}.
تهران: قلم، ۱۳۶۳.

\end{enumerate}

%============================================================
\section*{ب) منابع انگلیسی}

\begin{enumerate}[label={[\arabic*]}, leftmargin=2cm, resume]

\item
\en{Abrahamian, Ervand.
\textit{Khomeinism: Essays on the Islamic Republic}.
Berkeley: University of California Press, 1993.}

\item
\en{Abrahamian, Ervand.
\textit{A History of Modern Iran}.
Cambridge: Cambridge University Press, 2008.}

\item
\en{Abrahamian, Ervand.
\textit{Tortured Confessions: Prisons and Public Recantations
in Modern Iran}.
Berkeley: University of California Press, 1999.}

\item
\en{Abrahamian, Ervand.
\textit{Radical Islam: The Iranian Mojahedin}.
London: I.\,B.\,Tauris, 1989.}

\item
\en{Arjomand, Said Amir.
\textit{The Turban for the Crown: The Islamic Revolution in Iran}.
New York: Oxford University Press, 1988.}

\item
\en{Bakhash, Shaul.
\textit{The Reign of the Ayatollahs: Iran and the Islamic Revolution}.
New York: Basic Books, 1984.}

\item
\en{Beeman, William O.
\textit{The ``Great Satan'' vs.\ the ``Mad Mullahs'':
How the United States and Iran Demonize Each Other}.
Westport, CT: Praeger, 2005.}

\item
\en{Bowden, Mark.
\textit{Guests of the Ayatollah: The First Battle in
America's War with Militant Islam}.
New York: Atlantic Monthly Press, 2006.}

\item
\en{Brinton, Crane.
\textit{The Anatomy of Revolution}. Rev.\ ed.
New York: Vintage, 1965.}

\item
\en{Crist, David.
\textit{The Twilight War: The Secret History of
America's Thirty-Year Conflict with Iran}.
New York: Penguin, 2012.}

\item
\en{Dennett, Daniel C.
\textit{The Intentional Stance}.
Cambridge, MA: MIT Press, 1987.}

\item
\en{Elster, Jon.
\textit{Sour Grapes: Studies in the Subversion of Rationality}.
Cambridge: Cambridge University Press, 1983.}

\item
\en{Elster, Jon.
\textit{Ulysses and the Sirens: Studies in Rationality
and Irrationality}. Rev.\ ed.
Cambridge: Cambridge University Press, 1984.}

\item
\en{Festinger, Leon.
\textit{A Theory of Cognitive Dissonance}.
Stanford: Stanford University Press, 1957.}

\item
\en{Fingarette, Herbert.
\textit{Self-Deception}.
London: Routledge \& Kegan Paul, 1969.}

\item
\en{Franck, Thomas M.
\textit{The Power of Legitimacy Among Nations}.
New York: Oxford University Press, 1990.}

\item
\en{George, Alexander L., and Andrew Bennett.
\textit{Case Studies and Theory Development in the Social Sciences}.
Cambridge, MA: MIT Press, 2005.}

\item
\en{Howson, Colin, and Peter Urbach.
\textit{Scientific Reasoning: The Bayesian Approach}. 3rd ed.
Chicago: Open Court, 2006.}

\item
\en{Kant, Immanuel.
\textit{Groundwork of the Metaphysics of Morals} (1785).
Trans.\ Mary Gregor.
Cambridge: Cambridge University Press, 1998.}

\item
\en{Keddie, Nikki R.
\textit{Modern Iran: Roots and Results of Revolution}.
Updated ed.
New Haven: Yale University Press, 2006.}

\item
\en{Korsgaard, Christine M.
\textit{Creating the Kingdom of Ends}.
Cambridge: Cambridge University Press, 1996.}

\item
\en{Krippendorff, Klaus.
\textit{Content Analysis: An Introduction to Its Methodology}.
3rd ed.
Thousand Oaks, CA: Sage, 2013.}

\item
\en{Kunda, Ziva.
``The Case for Motivated Reasoning.''
\textit{Psychological Bulletin}, 108(3), 1990, pp.\,480–498.}

\item
\en{Kurzman, Charles.
\textit{The Unthinkable Revolution in Iran}.
Cambridge, MA: Harvard University Press, 2004.}

\item
\en{Lijphart, Arend.
``Comparative Politics and the Comparative Method.''
\textit{American Political Science Review}, 65(3), 1971, pp.\,682–693.}

\item
\en{MacIntyre, Alasdair.
\textit{After Virtue: A Study in Moral Theory}. 3rd ed.
Notre Dame, IN: University of Notre Dame Press, 2007.}

\item
\en{Martin, Vanessa.
\textit{Creating an Islamic State: Khomeini and the Making of
a New Iran}.
London: I.\,B.\,Tauris, 2000.}

\item
\en{Mele, Alfred R.
\textit{Self-Deception Unmasked}.
Princeton: Princeton University Press, 2001.}

\item
\en{Milani, Mohsen M.
\textit{The Making of Iran's Islamic Revolution:
From Monarchy to Islamic Republic}. 2nd ed.
Boulder, CO: Westview Press, 1994.}

\item
\en{Mill, John Stuart.
\textit{Utilitarianism} (1863).
Ed.\ Roger Crisp.
Oxford: Oxford University Press, 1998.}

\item
\en{Moin, Baqer.
\textit{Khomeini: Life of the Ayatollah}.
London: I.\,B.\,Tauris, 1999.}

\item
\en{Nickerson, Raymond S.
``Confirmation Bias: A Ubiquitous Phenomenon in Many Guises.''
\textit{Review of General Psychology}, 2(2), 1998, pp.\,175–220.}

\item
\en{Opwis, Felicitas.
\textit{Maṣlaḥa and the Purpose of the Law: Islamic Discourse
on Legal Change from the 4th/10th to 8th/14th Century}.
Leiden: Brill, 2010.}

\item
\en{Robertson, Geoffrey.
\textit{The Massacre of Political Prisoners in Iran, 1988:
A Report}.
Washington, DC: Abdorrahman Boroumand Foundation, 2010.}

\item
\en{Saltelli, Andrea, et al.
\textit{Global Sensitivity Analysis: The Primer}.
Chichester: Wiley, 2008.}

\item
\en{Schirazi, Asghar.
\textit{The Constitution of Iran: Politics and the State
in the Islamic Republic}.
London: I.\,B.\,Tauris, 1997.}

\item
\en{Sharot, Tali.
\textit{The Optimism Bias: A Tour of the Irrationally Positive Brain}.
New York: Pantheon, 2011.}

\item
\en{Skocpol, Theda.
\textit{States and Social Revolutions: A Comparative Analysis of
France, Russia, and China}.
Cambridge: Cambridge University Press, 1979.}

\item
\en{Trivers, Robert.
\textit{The Folly of Fools: The Logic of Deceit and
Self-Deception in Human Life}.
New York: Basic Books, 2011.}

\item
\en{Tufte, Edward R.
\textit{The Visual Display of Quantitative Information}. 2nd ed.
Cheshire, CT: Graphics Press, 2001.}

\item
\en{von Hippel, William, and Robert Trivers.
``The Evolution and Psychology of Self-Deception.''
\textit{Behavioral and Brain Sciences}, 34(1), 2011, pp.\,1–16.}

\end{enumerate}


%%%%%%%%%%%%%%%%%%%%%%%%%%%%%%%%%%%%%%%%%%%%%%%%%%%%%%%%%%%%
%%  INDEX — نمایه
%%%%%%%%%%%%%%%%%%%%%%%%%%%%%%%%%%%%%%%%%%%%%%%%%%%%%%%%%%%%

% Note: For full index, use makeidx or xindy.
% Below is a placeholder structure.
% In production, add \index{} commands throughout Segments 1-4
% and compile with: makeindex -s xindy.ist <jobname>

\cleardoublepage
\phantomsection
\addcontentsline{toc}{chapter}{نمایه}

% If using makeidx:
% \printindex

% Manual placeholder index:
\chapter*{نمایه}
\label{ch:index}

\begin{infoBox}[دربارهٔ نمایه]
نمایهٔ زیر مهم‌ترین اصطلاحات، اشخاص و رویدادها
را با ارجاع به شمارهٔ فصل فهرست می‌کند.
در نسخهٔ نهایی با استفاده از بستهٔ
\en{makeidx/xindy}
به‌صورت خودکار تولید خواهد شد.
\end{infoBox}

{%
\small
\begin{multicols}{2}
\subsection*{اشخاص}
\begin{description}[style=nextline, font=\normalfont]
\item[خمینی، روح‌الله] تمام فصل‌ها
\item[بازرگان، مهدی] ۶، ۷، ۱۶
\item[بختیار، شاپور] ۷
\item[بنی‌صدر، ابوالحسن] ۱۲، ۱۴، ۱۶
\item[بهشتی، محمد] ۸، ۱۴، ۱۶
\item[رفسنجانی، اکبر هاشمی] ۸، ۱۱، ۱۴، ۱۶
\item[سروش، عبدالکریم] ۴، ۱۴
\item[کدیور، محسن] ۴، ۸، ۱۳، ۱۴
\item[گنجی، اکبر] ۱۳
\item[منتظری، حسینعلی] ۱۲، ۱۶، پیوست ب
\item[یزدی، ابراهیم] ۶، ۷، ۱۶
\end{description}

\subsection*{مفاهیم و اصطلاحات}
\begin{description}[style=nextline, font=\normalfont]
\item[استدلال انگیزه‌محور] ۱۵
\item[ارزیابی بیزی] ۱۵
\item[تحلیل حساسیت] ۱۵
\item[تز ترکیبی مرحله‌ای] ۱۴، ۱۵، ۱۶
\item[تقیه] ۴، ۱۴
\item[خودفریبی] ۱۴، ۱۵، ۱۶
\item[ماتریس سازگاری] ۱۴
\item[مصلحت نظام] ۱۱، ۱۵، ۱۶، پیوست ب
\item[ناهماهنگی شناختی] ۱۵
\item[ولایت فقیه] ۴، ۸، ۱۱، ۱۳، پیوست ب
\item[ولایت مطلقهٔ فقیه] ۱۱، ۱۵، ۱۶، پیوست ب
\

<!-- POSSIBLE OVERLAP DETECTED (similarity: 85%) - REVIEW BELOW -->